% ================================================================
%  第 4 章 引言 (原书 p.109-122, 边码 P.I-P.XXII)
%  源文: _work/src1/task-ch04.md
% ================================================================
\chapter{引言}

% ---------------- p.109 ----------------
\section{(I) 五叶草}

\subsection{1. 梦与成就}

\origpage{I} 到今年七月,就是三年前了,我做了一个不大寻常的梦。我说“不大寻常”,这印象只是事后才有的,是醒来回想时才有的。梦本身来到我这里,却像世上最自然、最显然不过的事,不敲锣不打鼓——以至于醒来时我差点没留意它,随手把它扔进遗忘的角落,好去办“当天的正事”。从前一天起,我就已经动身去做一场关于我与数学之关系的反思了。这是我有生以来第一次肯费这个力气去看一看——而且,我之所以在那个时刻动手,实在是因为我几乎是被逼无奈、硬被赶上路的!在此前的那些月和那些年里,发生过一些如此奇特——若不是说暴烈——的事情,那是些数学激情的爆发,不打招呼就闯进我的生活,以至于实在再也没法继续不去看看到底发生了什么。

我说的这个梦没有任何情节,也没有任何动作。它只有一幅图像,静止不动,同时又极其鲜活。那是一个人的头,侧面像。看得出它正从右向左看。那是一个成年男子,无须,一头乱发在头的周围形成一圈力量的光轮。从这颗头上透出来的,尤其是那种年轻、欢快的力量感,仿佛从那柔韧有力的颈项弧线里涌出(那弧线是猜出来的,而不是看出来的)。脸上的表情,与其说是个成年人,或者一个已经沉稳下来的人(不管成熟与否),不如说是个顽皮的小淘气,正为某件他刚刚得手、或正在盘算着要干的勾当而洋洋得意。从那里透出来的,尤其是那种强烈、含而不露、迸发为嬉戏的生的欢乐……

% ---------------- p.110 ----------------
\origpage{II} 并没有第二个人在场,没有一个“我”在观看或凝视这另一个——我们只看得见他的头。但有的是对这头颅、对它透出的东西的一种强烈感知。也没有任何人去感受印象、去评论它、给它命名,或者给被感知的那个人贴上一个名字,称之为“某某”。有的只是这极其鲜活的东西,这颗男人的头,以及对这东西同样鲜活而强烈的感知。

至于醒来时,我并不是有意地,就记起了过去这一夜的梦;这个男人的头的影像,并没有以特别的强度从那堆梦里凸显出来,它并没有挤到前面来,冲我喊、向我低语:你该看的正是我!当这个梦出现在我对一夜之梦的匆匆一瞥里,在被窝里温暖的静谧中,我自然有了醒着的心智那个习惯性的反应:给看见的东西安上一个名字。何况我根本不必去找:只要我一问,我立刻就知道了,梦里那个在那儿出现的男人的头,不是别的,正是我自己的头。

这一手不赖,我当时想,到底还是得服气:在梦里这样看见自己,就像那是另一个人!这个梦的到来,有点像我在散步时,纯属最偶然的巧遇,撞上了一株四叶草,甚至五叶草,于是我按规矩惊异上片刻,然后继续赶路,就像什么也没发生过。

至少,事情差点就这样过去了。幸而——在这种情形下我碰过许多次——我还是为了问心无愧,把这个“不赖”的小插曲白纸黑字记了下来,并开始了一场反思,它本该顺着前一天的思路继续下去。可是,一针引一线,那一天的反思最终只是让我沉浸于这个不起眼的梦、这幅唯一的图像的含义,以及它带给我的关于我自己的那条消息。

这里不是细说我这一日的冥想教了我什么、带给我什么的地方。或者更确切地说,不是细说这个梦教了我什么、带给我什么的地方——那是在我已经进入一种专注与倾听的状态、从而能够接纳它要对我说的话之后。梦与这份倾听的第一枚直接的果实,是一股骤然涌来的新能量。正是这股能量支撑了随后几个月里那场长时间的冥想,与之相抗的是内心里顽强的阻力,我得靠耐心而执拗的功夫把它们一块一块拆掉。

我开始留意自己某些梦,已有五年;而这是第一个不以“信使之梦”如今已可辨认的那副面目出现的信使之梦——那种梦有惊人的场景手段,有异常强烈、有时令人翻肠倒肚的视象。这一个却“酷”到了家,没有任何强迫你注意的东西,简直是悄无声息——取亦可,舍亦可,毫无纠纷。

\origpage{III} 几周之前,倒来过一个老派风格的信使之梦,调门是戏剧性的,甚至是野性的;它突然而即刻地终结了一段漫长的数学狂热。两个梦之间唯一表面的亲缘关系是:两者之中都没有观察者。这个梦以一则简短有力的寓言,指出当时我生活里正在发生的一件事,而我却不肯费心去留意它——说实话,我还小心翼翼地不去理会它。正是这个梦让我当时明白了有一件反思的工作刻不容缓;几周后我便投入其中,它随后持续了将近六个月。\origpage{IV}在这部反思—见证《收获与播种》的最后一部分里,我多少有机会谈到它;这部作品开启本卷,并给了本卷它的名字\fnote{1}{}{尤其参见第43节《扫兴的老板——或压力锅》。}。

% ---------------- p.111 ----------------
\origpage{V} 我之所以用这另一个梦、用这幅我自己的影像—视象(我在德语笔记里称之为“Traumgesicht meiner selbst”)来开篇这则引言,是因为最近这几周里,这个梦的念头不止一次回到我这里,那时,那场“关于一位数学家的往昔”的冥想正走向它的终点。说实话,回头看,这个梦以来的三年在我眼里像是澄清与成熟的年月,走向对它那条单纯而清澈的消息的实现。梦把“我本来的样子”显示给我。同样清楚的是,在我醒着的生活里,我并不完全是梦所显示的那个人——来自久远的重负与僵硬(过去和现在都)常常妨碍我完全而单纯地做我自己。这些年里,尽管这个梦的念头只是偶尔回到我这里,它一定以某种方式起了作用。它绝不是某种范本或理想,让我努力去与之相像,而是对一种欢乐的单纯的悄声提醒——那单纯“就是我”,它以许多方式显现,并被召唤着去摆脱那些继续压着它的东西,去充分地绽放。这个梦是一条既纤细又有力的纽带,一头连着仍被来自过去的许多重负压着的现在,一头连着这现在已孕育在胚芽里的、近在眼前的“明天”——一个从此刻起就是我、而且肯定一直就在我之中的“明天”……

肯定,如果说最近这几周里这个难得提起的梦又一次鲜明地在场,那是因为在某个层面上——那不是探测与分析的思想的层面——我一定已经“知道”:我正在做、并正引向终点的这件工作,这件接续并深化了三年前那件工作的工作,是向着实现它带给我的、关于我自己的那条消息迈出的又一步。

如今在我看来,这就是《收获与播种》的主要意义所在,这件将近两个月的高强度工作的主要意义所在。直到它完成了,我才意识到我做这件事有多么重要。在这件工作的过程中,我经历过许多欢乐的时刻,那往往是一种促狭的、爱开玩笑的、兴高采烈的欢乐。也有过悲伤的时刻,有过一些时刻,我重新经历这些年里痛苦地击中过我的挫折与伤心——但没有过哪怕一刻的苦涩。我离开这件工作时,带着那种知道自己把一件工作做到了底的人的完全的满足。没有任何一件事,无论它多么“小”,是我在那里回避掉的,或者是我心里看重、本该说出而竟没有说出的,以至于此刻会在我心里留下一点不满、一点遗憾的残渣——无论那多么“小”。

在写这份见证时,我很清楚它不会讨所有人喜欢。甚至很可能,我找到了办法把所有人,一个不落,都得罪了。然而这根本不是我的本意,连得罪任何人也不是。我的本意不过是去看那些简单而重要的事,那些日常的事——我作为数学家的过去(有时也包括现在)里的事——好最终(晚知道总比不知道好!)毫无疑义、毫无保留地发现它们过去是什么、现在是什么;并且一路走一路用简单的话说出我所看见的。

\subsection{2. 一场旅行的精神}

这场反思最终成了《收获与播种》,它起初是作为《追寻栈》(À la Poursuite des Champs)第一卷(正在完成中)的“引言”开始的——《追寻栈》是我自1970年以来就打算发表的第一件数学工作。头几页是我去年六月在一个空闲的空当里写的;不到两个月前,我又从搁下的地方把这反思接了起来。我意识到有不少东西值得看、值得说,所以指望的是一篇比较充实的引言,三四十页。可是此后将近两个月里——直到此刻,我正在为这最初本是引言的文字写一篇新引言——我每天都以为今天就
% ---------------- p.112 ----------------
会把它做完,至不济明天或后天也就完了。等到几周之后我开始逼近一百页这个关口,这引言就被擢升为“导论性的一章”。又过了几周,当所谓“这一章”的规模远远超出正在准备的这卷里其他各章(写这几行时,除最后一章外都已写完),我终于明白了:它的位置不该在一本数学书里,这场反思与见证在那里实在太挤。它真正的位子是在单独的一卷里,那将是“数学反思”系列的第1卷——这个系列我打算今后几年顺着《追寻栈》的势头继续下去\fnote{2}{}{(5月28日)这里指的是《收获与播种》第一部分《愚妄与更新》的文本。写这几行时,第二部分尚未写出。}。

我不会说《收获与播种》——这部《数学反思》系列的第一卷(接下来先是《追寻栈》的两三卷)——是《反思》的“导论”卷。不如说,我把这第一卷看作将来之物的地基,或者更确切地说,看作定下基调的那一卷,给出我踏上这趟新旅程时所怀的那种精神的那一卷;这旅程我打算在今后几年里继续走下去,而它把我带向哪里,我说不上来。

关于本卷主干部分的这些说明,最后再交代几点实用性质的情况。读者在《收获与播种》的正文里偶见“本卷”的提法,不必奇怪——那暗指的是《追寻栈》的第一卷(《模型史》),我当时还以为自己仍在为它写引言。我不愿去“改正”这些地方,因为我最看重的是保住文本的自发性,保住它作为见证的真实性——不仅是对一个遥远过去的见证,也是对正在写作的当下这一刻的见证。

同样出于这个理由,我对文本初稿所作的修改,只限于纠正文笔上的笨拙,或某些有时含混、妨碍理解我要表达之意的说法。这些修改有时把我引向比写初稿时更清晰、更精细的领会。至于对初稿稍具实质性的改动——为给它添些层次、说得更精确、作些补充,或者(有时)加以更正——则放在大约五十条编号的“注”里,集中排在反思的末尾,它们占了全文四分之一还多。我用(1)之类的符号去指这些注……在这些注里,我挑出二十来条,觉得它们(以篇幅或实质而论)重要得可与这篇反思自发组织成的五十个“节”或“段”中的任何一节相比。这些较长的注也收进了目录,排在五十节的清单之后。不出所料,对某些长注来说,又需要给注再加一条或几条注。这些注便接在该注之后,用的是同一类指号;只有相当短的注例外,它们就排在同一页上作“脚注”,用的是诸如(46\textsubscript{1})、(46\textsubscript{2})那样的指号。

给正文的每一节、以及每一条最充实的注取一个名字,我从中得到很大的乐趣——何况后来这甚至成了我理清头绪的必需品。不用说,这些名字都是事后才找到的;当初开始写某一节或某条稍长的注时,我一条也说不出它的实质是什么。我用来称呼八个部分(I至VIII)的那些名字(如“劳动与发现”等等)就更是如此——这八部分是我事后把构成全文的五十节归并而成的。

至于这八个部分的内容,我只作极简短的评论。头两个部分——I(《劳动与发现》)与II(《梦与做梦者》)——含有一些关于数学工作、以及一般而论关于发现工作的反思要素。
% ---------------- p.113 ----------------
\origpage{VII}在其中,我本人牵涉得零散得多,也间接得多,远不如后面几个部分。后面这几个部分才尤其具有见证与冥想的品格。\origpage{VIII}III至VI部分,主要是关于我1948至1970年间“在数学世界里”那段数学家往昔的反思与见证。推动这场冥想的动机,首先是想理解这段过去,以此努力去理解并承担一个在某些方面有时令人失望或令人茫然的现在。VII(《孩子嬉戏》)与VIII(《孤独的历险》)两部分,更多涉及我从1970年至今、也就是从我离开“数学家的世界”不再回头以来,我与数学之关系的演变。那些动机、那些力量与境况,它们使我(连我自己都吃惊)在中断十三年多之后,重新开始一种“公开的”数学活动(即写作并发表《数学反思》)。
